%%
%%                  TEMPLATE for Math 404 project report
%%
\documentclass[11pt]{amsart}
%%% WARNING: Do NOT change the page size, fonts, or margins!  Penalties will apply.

\usepackage{verbatim}
\usepackage{graphicx}
\usepackage{amssymb,amsmath,amsthm}
\usepackage{placeins} %enables \FloatBarrier, which prevents figures and tables from going below it.
\usepackage{hyperref} %makes cross references into hyperlinks. 
\usepackage{caption}
\usepackage{subcaption} %Allows subfigures

\title{Modeling Climate Change Over Time}
\author{Japheth Carlson, Andy Criddle, Nathan Longhurst, Ethan Palenske}

% Change the date to match the date you actually wrote this paper
\date{\today} % or use \today


\begin{document}

\begin{abstract}
% Place abstract here. The abstract summarizes in one paragraph the main question, results, and conclusions drawn from your investigation.  
\par The goal of this paper is to examine historical climate change data and determine if we can find correlations between CO2 levels, changing temperatures, and shrinking sea ice levels. We also compare climate change with historical events, including volcanic eruptions and the popularization of cars, and try to determine future CO2, temperature, and ice trends.

\end{abstract}

\maketitle % this actually makes the title

%% First Section
\section{Problem Statement and Motivation}
% Give a clear statement of the problems or questions the project addresses,
% their context, and a compelling motivation for why they are worth studying.
% The problems or questions your project addresses should be original, meaningful,
% and reflect a deep understanding of the subject and its subtleties.  

% Make a clear case for why your question is interesting, well thought out,
% precisely formulated, and answerable, at least in principle, with adequate data
% and the techniques of machine learning.
% Also briefly review, with proper citations \cite{Vandermeersch}, what is already
% known about the research questions and what techniques others have used to study
% these questions.
% Explain the scope of the project and how it fits into existing research.

It is well established that Earth's atmospheric temperatures are seasonal.  
Aside from the four standard seasons associated with Earth's axial tilt, long-range seasonal temperatures are influenced by ocean temperatures and myriad minute details of Earth's orbit (called Milankovitch cycles). \cite{milankovitch1941canon}
While understanding these cycles is still an open area of investigation, they typically act in predictable ways over long periods of time.  
Within the last 200 years, however, Earth's climate has changed faster than in any time known since the mass extinction event of the dinosaurs \href{https://www.scientificamerican.com/article/todays-climate-change-proves-much-faster-than-changes-in-past-65-million-years/#:~:text=According%20to%20a%20Stanford%20University%20report%2C%20the,and%20global%20temperatures%20rise%205%20degrees%20Celsius}{(Source)}.

With the climate changes on Earth trending toward another potential mass extinction event, the immediate question becomes:  What has caused such a rapid change in global climate?  
This highlights the foundational importance of finding the correlation of Earth's surface/atmospheric temperature and other proposed correlating factors.  
Recently, there has been increasing momentum in both scientific and societal circles suggesting that Earth's rapid shift in temperature is largely due to a rise in atmospheric $CO^2$ levels. \cite{foote1856art, arrhenius1896xxxi, callendar1938artificial, keeling1960concentration, manabe1967thermal, charney1979carbon, santer1996search}
Moreover, concern has been raised about depleting sea ice temperatures, particularly in polar regions.  
We therefore begin our analysis by modeling these factors together over time, beginning as far back as data permits, until the present.

Following these initially intriguing questions, there is another core question of at what points our climate has shifted; while we can model the future using past measurements, those models will only work where they understand how that system has fundamentally changed along the way. 
This allows for more intelligent model selection for the ability to generalize into the future.



%%Second Section
\section{Data} 
% Discuss the sources, reliability, and suitability of your data for the problems you are addressing.  The dataset should be large enough and rich enough to give reliable, meaningful, and nuanced answers to the questions and problems addressed in your project. Remember that test-train-validation splits for sequential and time-series data are trickier than for tabular data. Be clear about how you handled splitting data into test, training and validation sets, and how you ensured that there was no leakage between test and training set and no other similar data problems.  Be clear about how you handled missing data.  Justify your choices on all these things. 

% \begin{figure} %Put your images in a figure like this 
% \begin{center}
%     \includegraphics[width=\textwidth]{Myfig.pdf}
% \end{center}
%      \caption[Behavior near $1/\sqrt{5}$]{Plots should be high resolution (pdf, 300 dpi), uncluttered, a reasonable size, and easily readable and understandable.  All figures should have a complete caption that helps the reader make sense of the figure even if they haven't read the paper yet. Here is an example: The generalized aliasing decomposition (GAD) decomposes risk (in black) as the sum of three parts: aliasing (blue) which generally peaks at the  interpolation threshold (when the model complexity matches the number of training points) and decreases to zero as the model complexity increases; model insufficiency (red), which dominates when models are too simple, and is monotonically decreasing as models become more complex; and data insufficiency (green), which is usually zero to the left of the interpolation threshold but increases monotonically as model complexity increases.}

% \label{fig:MeanSquaredError} % for automatic cross referencing
% \end{figure}

We have gathered data from various sources to analyze global trends and predict future temperatures. For average temperature data recorded throughout the globe, we have used data compiled by Berkeley Earth provided for free on \href{https://www.kaggle.com/datasets/berkeleyearth/climate-change-earth-surface-temperature-data/data}{Kaggle}. This dataset contains average temperatures recorded in various countries with both land and ocean measurements recorded monthly. This dataset dates back to 1750, but some earlier years/months were not recorded. Thus, we primarily examine the average land temperature measurements that date back to 1753 in our analysis to avoid issues with missing data. We note that in earlier recorded years, there is much more variation in the recorded average temperature. This may be due to the limited number of locations reporting temperature data.

To grasp the impacts of $CO_2$ and all greenhouse gas emissions on temperature, we use the emissions data provided by \href{https://ourworldindata.org/co2-and-greenhouse-gas-emissions}{Our World in Data} that contains most emissions caused by humans over the last 250+ years. 
To compare this to $CO_2$ emissions caused by volcanoes we use data collected by various experts in the area. \cite{fischer2019emissions}


%% Third Section
\section{Methods} 
% Your project should involve a thoughtful and original use of several of the
% machine learning methods and ideas we have learned in class and at least one
% idea or method we have not covered in class.  

% In this section you should describe and justify your selection of models,
% your choices of methods, any feature engineering you did, and the ways and
% reasons you chose your hyperparameters or network architectures.
% Your discussion should demonstrate a clear understanding of the principles
% involved and of the strengths and weaknesses of the models and methods used.  

% Japheth

% Andy

% Nathan

% Ethan

%%Fourth Section
\section{Results and Analysis} 
% Clearly and succinctly describe your results. Give a thorough analysis and a
% thoughtful discussion of the conclusions that can be drawn.
% Discuss the suitability and effectiveness of the different models and methods
% for the problems or questions treated.

% Japheth

% Andy

% Nathan

% Ethan


%%Fifth Section
\section{Ethical Implications and Conclusions}
% Thoughtfully analyze the ethical implications of your research questions,
% the data you gathered, and the analysis that was performed.
% Are there privacy or other implications from the collection or use of the data?
% Could your results and methods be misused or misunderstood?
% What can and should be done to prevent misuse and misunderstanding?
% Could your algorithms and methods result in a destructive self-fulfilling feedback loop?
% How could that be prevented or controlled?
% What other ethical implications does your work have?  




%% Sixth section
\section{Conclusion} 
% Describe the final conclusions that you draw from your computations, results, and analysis. 




%All text before this line should end before page 11 starts.
%%%%%%%%%%%%%%%%%%%%%%%%%%%%%%%%%%%%%
%% Bibliography below
%%%%%%%%%%%%%%%%%%%%%%%%%%%%%%%%%%%%%
\FloatBarrier % Keep all figures from being put after the bibliography
\newpage
%% If using bibtex, leave this uncommented
\bibliography{refs} %if using bibtex, call your bibtex file refs.bib
\bibliographystyle{alpha}

%% If not using bibtex, comment out the previous two lines and uncomment those below
%\begin{thebibliography}{99}
%\bibitem{Vandermeersch} First reference goes here
%\end{thebibliography}

\end{document}